\chapter*{Reading this proposal}
\addcontentsline{toc}{chapter}{Reading this proposal}
\label{sec:reading-guide}

The starting question is what a company can do after a shock when its published
accounts leave part of its history unknown. A fixed-rate borrower and a
floating-rate borrower can report the same interest expense today and need
different amounts of cash tomorrow. More journal detail helps only when it
preserves the contracts, dates, counterparties and decisions that determine
that response. Chapter~\ref{sim:observations} develops this example and the
joint family of histories behind it.

The first part follows the simulation problem from evidence to decisions.
Chapters~\ref{sim:reconstruction}--\ref{sim:scenarios} explain reconstruction,
learning and permitted information. Chapters~\ref{sim:debt} and
\ref{sim:equity} distinguish financing cash from contractual promises and
ownership. Chapters~\ref{sim:compression} and~\ref{sim:approximation} address
how to retain relevant uncertainty without storing a separate company at every
branch. The complete company example in Chapter~\ref{sim:company-case}, the
dated refinancing case in Chapter~\ref{sim:refinancing-case} and the solved
risky-debt model in Chapter~\ref{sim:finite-finance} make the consequences
concrete. Their supplied assumptions are part of each answer.

The second part turns those requirements into a proposed accounting system.
An analyst can start with the linked statements and actual disclosure in
Chapter~\ref{ch:case}, then follow the accounting chapters appropriate to the
company and reporting basis. An implementer should read the same case before
the language in Chapter~\ref{ch:language}, the interfaces in
Chapter~\ref{ch:implementation} and the operation reference in
Appendix~\ref{app:portable-interfaces}. The reference gives details that the
narrative motivates; it does not replace the economic and accounting premises
of an operation.

Basic double-entry accounting, financial statements and probability are assumed.
The text introduces the additional notation where it is used. Numerical cases
are teaching calculations unless an actual source is identified. A model can
be internally consistent while its chosen company description remains
incomplete; a successful calculation can also depend on an unverified contract
term or behavioral assumption. Chapter~\ref{ch:validation} separates the
checks needed for the specification, an implementation and predictive use.
The detailed contents that follow support targeted reference as well as a
sequential reading.
\clearpage
